\chapter{Pancreas}

\label{ch:ch47_pancreas}

The pancreas is a gland located deep in the abdomen, behind the stomach. It produces digestive enzymes and hormones such as insulin. Imaging of the pancreas is most commonly requested for two reasons: \textbf{acute pancreatitis} (sudden, severe upper abdominal pain with raised pancreatic enzymes) and \textbf{suspected pancreatic tumour} (painless jaundice, weight loss, or new-onset diabetes).

\textbf{CT is the primary imaging test for pancreas}. For acute pancreatitis, a contrast-enhanced CT is performed to assess the severity of inflammation and detect complications such as necrosis (dead pancreatic tissue) or fluid collections. For suspected pancreatic cancer, a dedicated \textbf{pancreas protocol CT} is performed---this uses thin cuts and a specific timing of contrast injection to maximise the visibility of the pancreatic gland and any masses.

\textbf{MRI} is used when CT is inconclusive, for characterising cystic pancreatic lesions, and for evaluating the pancreatic duct with MRCP (magnetic resonance cholangiopancreatography). \textbf{Endoscopic ultrasound (EUS)} provides the highest resolution images of the pancreas and allows biopsy (fine-needle aspiration) of suspicious lesions.

---

\section{Imaging}
Ultrasound may be limited by bowel gas. Pancreas-protocol CT, MRI with MRCP, and endoscopic ultrasound are selected for acute inflammation, duct disease, cystic lesions, or suspected neoplasm.

\begin{table}[h]
\centering
\caption{Pancreas imaging modalities}
\begin{tabular}{p{3cm} p{4cm} p{4cm}}
\hline
\textbf{Modality} & \textbf{Best for} & \textbf{Limitations} \\
\hline
CT (pancreas protocol) & Acute pancreatitis, pancreatic cancer staging, trauma & Radiation; contrast required; limited for small lesions \\
EUS + FNA & Tissue diagnosis of small lesions, cyst fluid analysis & Semi-invasive; operator-dependent; sedation needed \\
MRI/MRCP & Cystic lesion characterisation, small pancreatic lesions, duct evaluation & Longer acquisition; claustrophobia; less available \\
\hline
\end{tabular}
\end{table}

\begin{pearlbox}[Pancreas protocol CT]
A pancreas protocol CT uses thin collimation (1--2.5 mm), IV contrast with \textbf{pancreatic phase timing} (40--50 s post-injection), and portal venous phase (60--70 s). The pancreatic phase provides maximal enhancement of the pancreatic parenchyma and best demonstrates hypodense pancreatic adenocarcinoma against the brightly enhancing normal pancreas.
\end{pearlbox}

---

\section{Acute Pancreatitis}

\subsection{Diagnosis}
Acute pancreatitis is diagnosed clinically (revised Atlanta criteria):
\begin{itemize}
  \item Characteristic abdominal pain (acute onset, epigastric, often radiating to the back)
  \item Serum lipase $> 3 \times$ upper limit of normal
  \item Imaging confirmation (CT or MRI) is NOT required for diagnosis if clinical criteria are met
\end{itemize}

\textbf{CT is indicated when}:
\begin{itemize}
  \item Diagnostic uncertainty (atypical presentation)
  \item Clinical deterioration or failure to improve after 48--72 hours
  \item Suspected complications (necrosis, abscess, pseudocyst, vascular complication)
\end{itemize}

\subsection{CT Severity Index (Modified CT Severity Index --- MCTSI)}
The MCTSI is used to grade the severity of acute pancreatitis on CT:

\begin{table}[h]
\centering
\caption{Modified CT Severity Index (MCTSI)}
\begin{tabular}{p{4.5cm} p{3cm} p{3cm}}
\hline
\textbf{Finding} & \textbf{Points} & \textbf{Criteria} \\
\hline
$\leq 30\%$ pancreatic necrosis & 4 & Peripancreatic fluid collection without well-defined margins \\
$> 30\%$ pancreatic necrosis & 6 & Well-defined peripancreatic fluid collection or abscess \\
Extrapancreatic complications & 2 & Any of: pleural effusion, ascites, vascular complication (pseudoaneurysm, thrombosis), bowel involvement \\
Interstitial oedematous pancreatitis & 2 & Focal or diffuse pancreatic enlargement with peripancreatic fat stranding, but no necrosis \\
Normal pancreas & 0 & Normal enhancement, no peripancreatic abnormality \\
\hline
\multicolumn{3}{l}{\textbf{Severity}: 0--2 = mild; 4 = moderate; 6--10 = severe} \\
\hline
\end{tabular}
\end{table}

\begin{pearlbox}[CT severity correlates with outcomes]
Mild acute pancreatitis (MCTSI 0--2) has $< 5\%$ mortality and resolves with supportive care. Severe acute pancreatitis (MCTSI 6--10) carries 10--30\% mortality and requires ICU management. CT severity score is an independent predictor of pancreatic necrosis, organ failure, and death.
\end{pearlbox}

\subsection{Pancreatic Fluid Collections}
\begin{table}[h]
\centering
\caption{Pancreatic fluid collections (revised Atlanta classification)}
\begin{tabular}{p{3cm} p{3cm} p{3cm} p{3cm}}
\hline
\textbf{Collection} & \textbf{Timing} & \textbf{Type} & \textbf{Characteristics} \\
\hline
Acute necrotic collection (ANC) & $< 4$ weeks, necrotising & Heterogeneous, no well-defined wall & Contains necrotic debris \\
Acute peripancreatic fluid collection (APFC) & $< 4$ weeks, interstitial & Homogeneous, no wall & Resolves spontaneously \\
Pseudocyst & $> 4$ weeks, interstitial & Homogeneous, well-defined enhancing wall & Mature epithelialised wall \\
Walled-off necrosis (WON) & $> 4$ weeks, necrotising & Well-defined wall, heterogeneous & Contains necrotic debris; may require intervention \\
\hline
\end{tabular}
\end{table}

\subsection{Vascular Complications}
\begin{itemize}
  \item \textbf{Pseudoaneurysm}: Focal rounded enhancement adjacent to the pancreas (splenic, gastroduodenal, or pancreaticoduodenal arteries most common). Risk of rupture---requires angiographic embolisation.
  \item \textbf{Splenic vein thrombosis}: Results in left-sided (sinistral) portal hypertension with gastric varices.
  \item \textbf{Splanchnic vein thrombosis}: Portal or superior mesenteric vein thrombosis.
\end{itemize}

---

\section{Chronic Pancreatitis}
CT and MRCP findings of chronic pancreatitis:
\begin{itemize}
  \item \textbf{Pancreatic calcifications}: Pathognomonic; scattered throughout the parenchyma
  \item \textbf{Main pancreatic duct dilation}: $> 3$ mm with irregular caliber (``chain of lakes'' appearance on MRCP)
  \item \textbf{Parenchymal atrophy}: Small, atrophic gland
  \item \textbf{Ductal strictures}: Focal narrowing of the main duct
  \item \textbf{Pseudocysts}: Well-defined fluid collections within or adjacent to the pancreas
  \item \textbf{Vascular complications}: Pseudoaneurysms, venous thrombosis
\end{itemize}

\begin{pearlbox}[Autoimmune pancreatitis]
Autoimmune pancreatitis (AIP) is a steroid-responsive form of chronic pancreatitis that mimics pancreatic cancer. Key imaging features: diffuse ``sausage-shaped'' pancreatic enlargement with a hypodense peripheral halo (``capsule sign''); focal form may mimic a mass; no calcifications; associated with IgG4 elevation and other autoimmune conditions (e.g., retroperitoneal fibrosis, sialadenitis). Biopsy showing lymphoplasmacytic infiltrate with IgG4-positive cells confirms the diagnosis.
\end{pearlbox}

---

\section{Pancreatic Masses}

\subsection{Pancreatic Adenocarcinoma}
The most common pancreatic malignancy (85\% of pancreatic cancers). Arises from the pancreatic duct epithelium.

\textbf{CT findings}:
\begin{itemize}
  \item \textbf{Hypodense mass} in the pancreatic parenchyma on the pancreatic phase (most common location: head of pancreas)
  \item \textbf{Main pancreatic duct dilation} ($> 3$ mm) upstream of the mass (``double duct sign'' when both bile and pancreatic ducts are dilated)
  \item \textbf{Vascular involvement}: Abutment ($> 180^\circ$ encasement of SMA or celiac axis) or occlusion of SMV/portal vein---determines resectability
  \item \textbf{Atrophy} of the pancreatic parenchyma upstream of the obstruction
  \item \textbf{Metastases}: Liver, peritoneum, distant lymph nodes
\end{itemize}

\begin{table}[h]
\centering
\caption{Resectability criteria for pancreatic adenocarcinoma}
\begin{tabular}{p{3cm} p{4.5cm} p{4.5cm}}
\hline
\textbf{Category} & \textbf{SMV/portal vein} & \textbf{SMA / celiac axis} \\
\hline
Borderline resectable & $> 180^\circ$ or $\leq 180^\circ$ with contour abnormality or thrombosis & $\leq 180^\circ$ contact \\
Locally advanced & unresectable & $> 180^\circ$ contact with SMA or celiac axis \\
Metastatic & Distant metastases present & Distant metastases present \\
Resectable & No contact or $\leq 180^\circ$ without contour abnormality & No contact \\
\hline
\end{tabular}
\end{table}

\subsection{Pancreatic Neuroendocrine Tumours (pNETs)}
\begin{itemize}
  \item \textbf{CT}: Hypervascular mass (enhances intensely in the arterial phase), unlike the hypovascular adenocarcinoma
  \item \textbf{Functional pNETs}: Insulinoma (most common functional; small, hypervascular, often $< 2$ cm), gastrinoma (associated with Zollinger--Ellison syndrome), glucagonoma, VIPoma
  \item \textbf{Non-functional pNETs}: Often larger at diagnosis; incidentally discovered
  \item \textbf{Grading}: Ki-67 index and mitotic rate guide behaviour (G1, G2, G3)
\end{itemize}

\begin{pearlbox}[Hypervascular vs hypovascular pancreatic mass]
A \textbf{hypervascular} pancreatic mass in the arterial phase is typical of a neuroendocrine tumour. A \textbf{hypovascular} (hypodense) mass in the pancreatic phase is typical of adenocarcinoma. This single distinction is fundamental to the differential diagnosis.
\end{pearlbox}

\subsection{Cystic Pancreatic Lesions}

\begin{table}[h]
\centering
\caption{Cystic pancreatic lesions}
\begin{tabular}{p{2.5cm} p{3cm} p{3.5cm} p{3cm}}
\hline
\textbf{Lesion} & \textbf{CT/MRI features} & \textbf{MRCP} & \textbf{Management} \\
\hline
IPMN (branch duct) & Side-branch cyst $> 5$ mm; adjacent to main duct & Communication with side branch & Based on high-risk stigmata (size $> 3$ cm, mural nodules, main duct dilation) \\
IPMN (main duct) & Dilated main pancreatic duct $> 5$ mm & Communication with main duct & High malignant potential; surveillance or surgery \\
Mucinous cystic neoplasm (MCN) & Macroscopic, thick-walled, mural nodules; body/tail of pancreas; middle-aged women & No duct communication & Surgical resection (malignant potential) \\
Serous cystadenoma & Microcystic (honeycomb), central scar with calcification & Stellate septations & Usually benign; surveillance \\
Solid pseudopapillary tumour & Solid-cystic mass with haemorrhage; young women & Variable & Surgical resection \\
\hline
\end{tabular}
\end{table}

\begin{pearlbox}[Sendai/Fukuoka guidelines for IPMN]
High-risk stigmata requiring surgery: main duct IPMN, mural nodules $> 5$ mm, main duct dilation $> 10$ mm, cytology suspicious for malignancy. Worrisome features warranting EUS: cyst $> 3$ cm, thickened walls, non-enhancing mural nodules, abrupt main duct change, lymphadenopathy.
\end{pearlbox}

---

\section{When to Recommend EUS/FNA}
EUS with fine-needle aspiration (FNA) is recommended when:
\begin{itemize}
  \item A cystic lesion cannot be confidently characterised on CT/MRI
  \item Cyst fluid analysis is needed (CEA level, amylase, cytology)
  \item A solid pancreatic mass is not accessible by percutaneous biopsy
  \item Focal autoimmune pancreatitis needs to be distinguished from cancer
  \item Mural nodules in an IPMN need histological assessment
\end{itemize}

---

\section{Interpretation}
Assess gland size and enhancement, peripancreatic inflammation, fluid collections, necrosis, ducts, vessels, adjacent organs, and distant disease. Describe cyst morphology, mural nodules, duct communication, and interval change when relevant.

\section{Safety}
Consider contrast, renal function, pregnancy, sedation, and the urgency of suspected infection or ischemia. Report necrosis, hemorrhage, vascular complication, obstruction, and suspicious solid components urgently.

\clearpage
